%BeginMC
\question
Violation of these two Hardy-Weinberg assumptions always causes homozygosity to increase.

\begin{choices}
\correctchoice Genetic drift and non-random mating.%EndChoice
\choice Mutation and natural selection.%EndChoice
\choice Natural selection and gene flow.%EndChoice
\choice Genetic drift and natural selection.%EndChoice
\choice Non-random mating and gene flow.%EndChoice
\choice Gene flow and mutation.%EndChoice
\end{choices}
%EndMC


%BeginMC
\question
The biological species concept emphasizes

\begin{choices}
\correctchoice reproductive isolation between species.%EndChoice
\choice the formation of cryptic species.%EndChoice
\choice gene flow between species.%EndChoice
\choice reproductive isolation between populations.%EndChoice
\choice the formation of morphologically distinct species.%EndChoice
\end{choices}
%EndMC

%BeginMC
\question
Beetle pollinators of a particular plant are attracted to the bright orange flowers. The beetles not only pollinate the flowers, but they mate while inside of the flowers. A mutant version of the plant with red flowers becomes more common over time. A particular variant of the beetle prefers the red flowers to the orange flowers. Over time, these two beetle variants diverge from each other to such an extent that interbreeding is no longer possible. What kind of speciation has occurred in this example, and what has driven it?

\begin{choices}
\choice allopatric speciation; habitat isolation%EndChoice
\correctchoice sympatric speciation; habitat isolation%EndChoice
\choice allopatric speciation; temporal isolation%EndChoice
\choice sympatric speciation; temporal isolation%EndChoice
\choice allopatric speciation; behavioral isolation%EndChoice
\choice sympatric speciation; gametic isolation%EndChoice
\end{choices}
%EndMC

%BeginMC
\question
Two species of \textit{Heliconius} butterflies occupy different habitats but the two species occur together in one area where their habitats are most similar. In this area, the two species often interbreed and hybrid offspring have been observed.  Yet, biologists continue to classify the butterflies as two separate species.  Which of the following observations would provide the best support for classifying the butterflies as two separate species? 

\begin{choices}
\choice The two species are reproductively isolated because they occupy different habitats.%EndChoice
\choice The two species are cryptic species.%EndChoice
\choice The hybrids are cryptic species.%EndChoice
\correctchoice The hybrids can't reproduce.%EndChoice
\choice The two species are morphological species, not biological species.%EndChoice
\end{choices}
%EndMC

%BeginMC
\question
Over time, the movement of people among all countries on Earth has steadily increased. This has altered the course of human evolution by increasing

\begin{choices}
\choice non-random mating. %EndChoice
\choice geographic isolation. %EndChoice
\choice genetic drift. %EndChoice
\correctchoice gene flow. %EndChoice
\choice natural selection. %EndChoice
\end{choices}
%EndMC

%BeginMC
\question
Which of these conditions should completely prevent the occurrence of natural selection in a population over time?

\begin{choices}
\correctchoice All variation between individuals is due only to environmental factors. %EndChoice
\choice The environment is changing at a relatively slow rate. %EndChoice
\choice The population size is large. %EndChoice
\choice The population lives in a habitat where there are no competing species present. %EndChoice
\choice The population mates non-randomly.%EndChoice
\end{choices}
%EndMC

%BeginMC
\question
A species of sport fish reaches an average adult length of 100 cm. Fishing regulations require that only individuals longer than 75 cm may be kept. Shorter fish must be released because they have greater reproductive output relative to large individuals. Based on your knowledge of natural selection, you should predict that the average length of the adult fish population will

\begin{choices}
\correctchoice remain unchanged. %EndChoice
\choice gradually decline. %EndChoice
\choice rapidly decline. %EndChoice
\choice gradually increase. %EndChoice
\choice rapidly increase. %EndChoice
\end{choices}
%EndMC

%BeginMC
\question
\Krakatou is a volcanic island in Indonesia.  An eruption in 1883 killed all living organisms on the island.  Many of organismal populations that are present now tend to show low genetic variation.  This is likely due to

\begin{choices}
\choice directional selection that favored traits that allow organisms to withstand volcanic activity. %EndChoice
\choice balancing selection that favored traits that allow organisms to withstand volcanic activity. %EndChoice
\choice a genetic bottleneck caused by the eruption, followed by regrowth of the population. %EndChoice
\correctchoice a series of founder effects with different species recolonizing from surrounding islands. %EndChoice
\choice The small population size causing random fluctuation of allele frequencies. %EndChoice
\end{choices}
%EndMC


%BeginMC
\question
According to Darwin's concept of descent with modification:

\begin{choices}
\choice Populations evolve by natural selection acting on individuals. %EndChoice
\correctchoice Species evolve from ancestral forms through the accumulation of different adaptations. %EndChoice
\choice Individuals with greater fitness will evolve more quickly than individuals with lower fitness. %EndChoice
\choice All species have evolved from a single common ancestor. %EndChoice
\choice All organisms overproduce offspring, leading to differential reproductive success. %EndChoice
\choice All organisms overproduce offspring, leading to different adaptations. %EndChoice
\end{choices}
%EndMC

%BeginMC
\question
Two closely related species of maple trees occur together in the southeastern United States.  Pollen from each species is distributed by the wind.  If the pollen from one species lands on the stigma (female pollen receptor) of the other species but a protein interaction between the pollen and stigma prevents fertilization, then which reproductive barrier keeps the two maple species isolated?

\begin{choices}
\choice mechanical.%EndChoice
\choice behavioral.%EndChoice
\choice habitat.%EndChoice
\correctchoice gametic.%EndChoice
\choice temporal.%EndChoice
\end{choices}
%EndMC


%BeginMC
\question
In a population with two alleles, \textit{A} and \textit{a}, that are in Hardy-Weinberg equilibrium, the frequency of allele \textit{a} is 0.2. What is the frequency of heterozygotes? 

\begin{choices}
\choice 0.4 %EndChoice
\choice 0.02 %EndChoice
\choice 0.04 %EndChoice
\choice 0.16 %EndChoice
\correctchoice 0.32 %EndChoice
\end{choices}
%EndMC

%BeginMC
\question
The rabbitsfoot mussel was once widepsread throughout rivers of eastern North America but now remains in only a few very small populations. The decrease occurred rapidly due to habitat loss. Genetics studies show that the remaining populations have little or no genetic variation. The best explanation for this result is:

\begin{choices}
\choice Balancing selection favored only one or two alleles. %EndChoice
\correctchoice Genetic drift caused the loss of heterozygosity. %EndChoice
\choice Inbreeding caused an increase in heterozygosity. %EndChoice
\choice Directional selection favored only one or two alleles. %EndChoice
\choice A bottleneck caused the loss of homozygosity. %EndChoice
\end{choices}
%EndMC

%BeginMC
\question
The overuse of drugs such as penicillin has caused an increase in the number of bacterial species that are resistant to the drugs. The most likely explanation for this result is:

\begin{choices}
\choice Most bacteria can resist drugs. %EndChoice
\choice Penicillin was not an effective drug. %EndChoice
\choice Bacteria began making proteins resistant to the drugs. %EndChoice
\choice The drugs caused mutations in the bacteria that gave them resistance. %EndChoice
\correctchoice A few bacteria were naturally resistant and natural selection favored their increase. %EndChoice
\end{choices}
%EndMC

%BeginMC
\question
Which mouse has the greatest relative fitness? (A clutch is the result of one breeding event.)

\begin{choices}
\choice One that lives for 4 years, has 3 clutches per year, and 4 offspring per clutch.%EndChoice
\choice One that lives for 3 years, has 3 clutches per year, and 5 offspring per clutch.%EndChoice 
\choice One that lives for 3 years, has 2 clutches per year, and 6 offspring per clutch. %EndChoice
\choice One that lives for 2 years, has 3 clutches per year, and 8 offspring per clutch. %EndChoice
\correctchoice One that lives for 3 years, has 3 clutches per year, and 7 offspring per clutch. %EndChoice
\end{choices}
%EndMC

%BeginMC
\question
(Think carefully about this.) Organisms living in the northern rocky intertidal face competing selective forces. Larger individuals benefit by having a larger internal volume, which helps to conserve body temperature. However, smaller individuals are less likely to be dislodged by large waves from storms. Thus, selection may favor intermediate body size for any particular species. This is an example of

\begin{choices}
\choice directional selection. %EndChoice
\choice stabilizing selection. %EndChoice
\choice disruptive selection. %EndChoice
\correctchoice balancing selection. %EndChoice
\choice natural selection. %EndChoice
\end{choices}
%EndMC
